\chapter{Équations, inéquations et signes}
\section*{Exercice 6.1 -- Premier degré}
La première donne $3x=15$, donc $S=\{5\}$. La seconde conduit à $-3=2$, impossible : $S=\varnothing$.

\section*{Exercice 6.2 -- Produit nul}
Le produit s'annule si $2x-1=0$ ou $x+4=0$, donc pour $x=1/2$ ou $x=-4$. Pour le carré, $x=\pm\sqrt{12}=\pm2\sqrt3$.

\section*{Exercice 6.3 -- Quotient}
Dans les deux cas, $x\ne1$. La première donne $x+3=2x-2$, donc $x=5$, admissible. La seconde se simplifie en $x+1=2$ sur le domaine, ce qui donne le seul candidat $1$, interdit. Son ensemble de solutions est vide.

\section*{Exercice 6.4 -- Signes}
Le produit est positif à l'extérieur des racines : $]-\infty,-2[\cup]3,+\infty[$. Le quotient est négatif entre $-2$ et $3$, nul en $3$ et interdit en $-2$ : solution $]-2,3]$.

\section*{Exercice 6.5 -- Budget}
Pour $n\in\N$, $14+3n\leq50$ équivaut à $n\leq12$. Le maximum est $12$ séances, pour un coût exactement égal à $50$ euros.

\section*{Exercice 6.6 -- Équivalence ?}
Il a supposé implicitement $x\ne0$ ; la division n'est donc pas équivalente sur tout le domaine. Le produit nul donne les deux solutions $0$ et $2$. Une autre méthode consiste à traiter d'abord $x=0$, puis le cas $x\ne0$.
